\section{Discussion}
\label{sec:discussion}

\subsection{What same-event evaluation misses}
On Sen1Floods11's own split HAND looks useless (+0.007); on an event outside training (Bolivia) and on 79
external floods it helps clearly. The contrast is what a practitioner would see when choosing inputs from
the standard benchmark. It does not fully isolate the effect of the split: the external events also differ
from Sen1Floods11 in region, scenes and labelling, so dataset shift and the change from known to new events
act together. Bolivia, which holds the dataset fixed, points the same way but is a single event; a
controlled within-dataset comparison of same-event and held-out-event splits at matched budgets would
separate the two.

\subsection{What makes HAND useful}
The model uses what HAND says about the place. The event-level gain disappears when HAND is shifted by
500\,m in each axis, blurred or taken from another flood, and a model trained on mismatched HAND gains
nothing; DEM and slope recover only a quarter of it. HAND is height relative to the local drainage
network, which relief and slope only approximate, and its position to within a few hundred metres and its
detail at that scale matter. Where it helps is consistent with this: HAND removes false alarms on higher
ground and on steep, shadowed slopes. Pooled over pixels, part of the gain survives a 500-m shift and DEM
and slope recover more of it, plausibly because the largest floods cover broad low-lying plains that stay low after
a shift and that relief alone identifies; the event-level gain, which weights every flood equally, depends
on the finer drainage structure. Terrain priors say where water \emph{can} be; a permanent-water map says
where it \emph{has been}, and with 159 training events it appears to teach the model to look for water only there
(pooled flood recall 0.69 → 0.57 at unchanged precision). As a physics loss, by contrast, HAND adds nothing: trained models already
satisfy the tested constraints (monotonicity is violated on fewer than 0.5\% of pixel pairs), and enforcing
them removes violations that are not where the errors are. Here the prior helps as an input, which
carries information the SAR does not.

\subsection{When the advantage shrinks or fails}
\emph{Many training events.} A model that has seen few floods may not have learned which backscatter
patterns are terrain artefacts, and HAND tells it directly; a model that has seen many floods could learn
this from SAR alone. The endpoints fit this hypothesis---a clear gain with 11 Sen1Floods11 events, none
detectable on flood extent with 159 GEOID-Flood events---but they differ in dataset too, and the
pre-registered test of a decline in between was null; draw-to-draw variance at 20--40 events is large. With
20 GEOID-Flood training events, DEM and slope also recover most of HAND's gain, which suggests that more
varied training lets a model infer where water collects from relief; we did not test this directly.
\emph{The pre-event image} partly substitutes for terrain with few events (pooled gain at 20 events +0.024, not significant, vs.\ +0.066).

\emph{Failure cases.} HAND fails on flat terrain where everything is near drainage (Somalia, EMSR419), on
water bodies the DEM drainage network misses (the Swedish lakes), and in mangroves (pooled over the seven
events that contain them). Neither pre-registered remedy works: gating HAND on flat tiles removes gains, and HAND dropout
shrinks the model toward SAR only. The losses sit in a few events that terrain statistics do not pick out
in advance; a fix would need to know when the DEM's drainage network is wrong. The stratified results are
descriptive: land-cover classes are tied to geography and flood type, so they show where the gain occurs,
not why, and the absolute IoU in built-up areas stays low even with HAND.

\emph{Registration.} Because the gain depends on HAND's position, terrain and radar must be co-registered.
Sentinel-1 IW GRD products are specified to about 7\,m absolute location accuracy~\cite{s1iwgrd}, and HAND
cells are 30\,m; both are far below the 500-m shift that removes the event-level gain, and a 100-m shift
costs little. Resampling HAND onto each SAR tile's own grid, as here, avoids larger errors.

\subsection{Limitations}
Besides the caveats stated with each result (one TerraMind configuration; a fixed step budget; 3--5 seeds),
GEOID-Flood labels are CEMS rapid-mapping products~\cite{cems}, not hand labels; their quality varies by
activation, and we did not audit them, so some apparent false alarms may be unmapped water. 50 of the 79
test events are in Europe. The change-detection baseline stacks the pre-event image as two extra channels.
HAND inherits the drainage-network errors of the GLO-30 DEM, and we did not compare it with other
hydrologically derived terrain features such as distance to drainage, or with HAND from other DEMs.
